% from ega0-19-fr.tex
\oldpage[0\textsubscript{IV-1}]{102}

élément de $A'$ étant inversible dans $A$, le corps des fractions de $A'$
(isomorphe à $\mathbf Q$) est aussi contenu dans $A$.

Prouvons enfin que $a)$ entraîne $c)$~; la condition $a)$ entraîne que $A$
contient un sous-corps premier $k_0$, isomorphe au sous-corps premier de $k$,
et dont $k$ est par suite extension séparable. Appliquant (19.6.2), on voit
d'abord qu'il existe un isomorphisme $f$ de $k$ sur un sous-corps $K$ de $A$.
D'autre part, soit $(x_i)_{1\leq i\leq n}$ un système d'éléments de
$\mathfrak m$ tel que les classes mod.~$\mathfrak m^2$ des $x_i$ forment une
base (sur $k$) de $\mathfrak m/\mathfrak m^2$. Puisque $A$ est complet, il y a
alors un homomorphisme continu d'anneaux
$u:k[[T_1,\ldots,T_n]]\to A$ tel que $u$ soit égal à $f$ dans $k$ et que
$u(T_i)=x_i$ pour tout $i$, et cet homomorphisme est surjectif en vertu du
choix des $x_i$ (Bourbaki, \emph{Alg. comm.}, chap.~III, \textsection2,
n\textsuperscript{o}~9, prop.~11).
\end{proof}

\begin{theorem}[19.6.4]
\label{0.19.6.4-fr}
Soient $k$ un corps, $A$ une $k$-algèbre locale noethérienne, $\mathfrak m$ son
idéal maximal, $K=A/\mathfrak m$ son corps résiduel, $A$ étant muni de la
topologie $\mathfrak m$-préadique. On suppose que $K$ soit une extension
séparable de $k$. Alors les conditions suivantes sont équivalentes~:
\begin{enumerate}
  \item[a)] $A$ est une $k$-algèbre formellement lisse.
  \item[b)] Le complété $\widehat A$ de $A$ est $k$-isomorphe à un anneau de
  séries formelles $K[[T_1,\ldots,T_n]]$ (dont la structure de $k$-algèbre est
  définie par l'homomorphisme composé
  \[
    k\xrightarrow{\varphi}K\xrightarrow{\psi}K[[T_1,\ldots,T_n]],
  \]
  où $\psi$ est l'injection canonique et $\varphi$ l'homomorphisme définissant
  la structure d'extension de $K$).
  \item[b')] $\widehat A$ est un anneau isomorphe à un anneau de séries
  formelles $K[[T_1,\ldots,T_n]]$.
  \item[c)] $A$ est un anneau local régulier.
\end{enumerate}
\end{theorem}

\begin{proof}
Le fait que $a)$ entraîne $b)$ est le cas particulier de ((19.5.4),
équivalence de $a)$ et $c)$), appliqué en y remplaçant $A$ par $k$, $B$ par
$A$ et $\mathfrak J$ par $\mathfrak m$, compte tenu de (19.6.1). Il est clair
que $b)$ entraîne $b')$ et $b')$ entraîne $c)$ par (17.1.1). Enfin si $c)$ est
vérifiée, il résulte de (17.1.1, $a)$) et de ((19.5.4), équivalence de $b)$ et
$a)$), appliqué à $K=A/\mathfrak m$ remplaçant $C$, que $A$ est une
$k$-algèbre formellement lisse.
\end{proof}

\begin{corollary}[19.6.5]
\label{0.19.6.5-fr}
Soient $k$ un corps, $A$ une $k$-algèbre locale noethérienne, $\mathfrak m$ son
idéal maximal, $A$ étant muni de la topologie $\mathfrak m$-préadique.
Supposons que $A$ soit une $k$-algèbre formellement lisse~; alors $A$ est
géométriquement régulier sur $k$, autrement dit (cf.~$\mathbf{IV}$, 6.7.6),
pour toute extension finie $k'$ de $k$, l'anneau semi-local
$A'=A\otimes_k k'$ est régulier (17.3.6). En outre, si $K$ est le corps
résiduel de $A$, $\widehat A$ est isomorphe à un anneau de séries formelles
$K[[T_1,\ldots,T_n]]$.
\end{corollary}

\begin{proof}
Comme $\widehat A'=\widehat A\otimes_k k'$, il résulte de (19.3.6) et de
(17.1.5) (appliqué aux anneaux locaux en les idéaux maximaux de $A'$) qu'on
peut se borner au cas où $A$ est complet~; alors $A'$ est une $k'$-algèbre
formellement lisse (19.3.5, (iii)) et par ailleurs est composée directe de
$k'$-algèbres locales, qui sont aussi formellement lisses (19.3.5, (v)). On est
donc ramené à prouver que $A$ est \emph{régulier}. Soit $k_0$ le sous-corps
premier de $k$~; comme $k$ est une extension séparable de $k_0$, c'est une
$k_0$-algèbre formellement lisse (19.6.1) et en vertu de l'hypothèse, il en est
de même de $A$ (19.3.5, (ii)). Comme le corps résiduel $K$ de $A$ est une
extension séparable de $k_0$, on peut alors appliquer (19.6.4) à $A$ et $k_0$,
d'où la conclusion.
\end{proof}

